% 麦克斯韦妖（综述）
% license CCBYSA3
% type Sum

本文根据 CC-BY-SA 协议转载翻译自维基百科\href{https://en.wikipedia.org/wiki/Maxwell\%27s_demon}{相关文章}

麦克斯韦妖（Maxwell's demon）是一个思想实验，看似可以推翻热力学第二定律。它由物理学家詹姆斯·克拉克·麦克斯韦于 1867 年提出。[1] 在他最初的信件中，麦克斯韦称这个存在为“有限的生命”或“能够与分子玩一场技巧游戏的存在”。后来开尔文勋爵称其为“妖”。[2]

在这个思想实验中，一个妖怪掌管着连接两间气体腔室的门。当单个气体分子（或原子）接近这扇门时，妖怪会迅速开关门：只允许高速运动的分子往一个方向通过，只允许低速运动的分子往另一个方向通过。由于气体的动能温度依赖于其分子的速度，妖怪的操作会使一间腔室升温，另一间腔室降温。这将导致系统的总熵减少，而似乎并未施加任何功，从而违反了热力学第二定律。

麦克斯韦妖的概念在科学哲学和理论物理学中引发了大量争论，直至今日仍在继续。它也推动了人们对热力学与信息论关系的研究。大多数科学家基于理论理由认为，没有任何装置能够以这种方式违反第二定律。也有一些研究者在实验中实现了某种形式的麦克斯韦妖，但它们都在某种程度上与思想实验不同，且没有一个被证明真正违背了热力学第二定律。
\subsection{这一思想的起源与发展}
这个思想实验最早出现在 1867 年 12 月 11 日麦克斯韦写给彼得·格思里·泰特的信中。1871 年，他又在写给约翰·威廉·斯特拉特（John William Strutt，后来的瑞利勋爵）的信里再次提到这一点，随后在 1872 年出版的热力学著作《热论》中正式向公众呈现。[3]

在信件和书籍中，麦克斯韦将那个掌管两腔室之间之门的存在描述为一个“有限的生命”。由于他是一个虔诚的宗教信徒，他从未使用过“demon（恶魔）”一词。相反，威廉·汤姆孙（William Thomson，即开尔文勋爵）才是第一个在 1874 年《自然》期刊中将麦克斯韦的概念称作“demon”的人。他暗示自己所指的是希腊神话中的 daemon，即在幕后起作用的超自然存在，而不是一个邪恶的存在。[2][4][5]
\subsection{原始思想实验}
热力学第二定律保证了（通过统计概率），当两个温度不同的物体彼此接触并与宇宙其他部分隔离时，它们最终会演化到一种热力学平衡状态，即二者的温度大致相同。[6] 热力学第二定律也可以表述为：在一个孤立系统中，熵永不减少。[6]

麦克斯韦构思了一个思想实验，以加深对第二定律的理解。他对实验的描述如下：[6][7]

……假设我们设想一种存在，它的能力如此敏锐，以至于能够追踪到每一个分子的运动轨迹。这样一种存在，其属性和我们一样本质上是有限的，却能够做到我们不可能做到的事情。因为我们已经看到，在一个充满均匀温度空气的容器中，分子的速度并不完全相同，尽管任意选取的大量分子的平均速度几乎完全一致。现在，假设这样的容器被隔板分为两部分，A 和 B，隔板上有一个小孔。有一个能够看到单个分子的存在，负责开关这个小孔，使得只有速度较快的分子能够从 A 进入 B，而只有速度较慢的分子能够从 B 进入 A。这样，他无需做功，就能使 B 的温度升高、A 的温度降低，从而与热力学第二定律相矛盾。

换句话说，麦克斯韦设想了一个被分隔为两部分 A 和 B 的容器。[6][8] 两部分都充满了相同的气体，并且初始温度相同，彼此相邻。一个假想的妖怪观察着两边的分子，并把守着两部分之间的活门。当 A 中一个快于平均速度的分子飞向活门时，妖怪便将其打开，让该分子从 A 飞入 B。同样地，当 B 中一个慢于平均速度的分子飞向活门时，妖怪会让它从 B 进入 A。结果是，B 部分分子的平均速度提高，而 A 部分的平均速度降低。由于平均分子速度对应于温度，因此 A 的温度下降，B 的温度上升，这与热力学第二定律相悖。此时，若在热源 A 和 B 之间运转一台热机，就能够从这种温差中提取有用功。

值得注意的是，妖怪必须允许分子双向通过，才能仅仅产生温度差；如果它仅允许快于平均速度的分子单向从 A 进入 B，那么 B 一侧不仅会升温，还会产生更高的压力。
